\section{Introduction} \label{introduction}

In the \textit{Online Graph Exploration} problem, an agent is initially
placed onto a random node of an unknown but labeled weighted graph. The agents task is to visit
all nodes of the graph and return to the starting node.
Since the agent gets the information about the graph online, it only knows about the already visited nodes,
as well as the neighborhood of these nodes at any point in time. Whenever a new node is visited,
the adjacent edges and the neighborhood of the visited node are revealed to the agent.
Traversing an edge incurs the cost of its weight. The goal of the agent is to explore the graph
as efficiently as possible, trying to keep down the cost of the exploration.
This problem was initially described as \textit{the online traveling salesperson problem} 
by Kalyanasundaram and Pruhs \cite{Kalyanasundaram1993}.\\
We examine a modified version of this problem using multiple agents starting at the same node, which
have to collaboratively explore a graph so that each node is visited at least once by some agent.
This variation does not have a consistent name in the literature and will here be referred to as 
\textit{Multi-Agent Online Graph Exploration}.

\subsection{Motivation}

The Online Graph Exploration problem has been widely used to describe the exploration of unknown terrain by multiple autonomous robots like mars rovers \cite{Dynia}, or robot vacuums \cite{Ortolf2012},
where the explored terrain is modeled as an unknown graph. It is also used to describe fleets of robots fully exploring unknown caves or damaged structures, to find and rescue human beings \cite{Higashikawa2012}. 
The literature contains multiple variations of this problem in terms of cost and communication models. Some allow agents to freely
communicate, while in some variations the agents have to meet to exchange information, or exchange information by writing and reading to/from the nodes they are currently located at. Some variations focus on the general exploration time (time model), while others 
focus on the maximum distance traversed by an individual agent. This model is used to describe the maximum energy consumption of any single agent during an exploration (energy model).\\
In the single-agent case, many graph classes like directed graphs \cite{Foerster2016}, tadpole graphs \cite{Brandt2020} or cactus graphs \cite{Fritsch2021} have already been investigated, 
while the main focus of the literature in the case of multi-agent exploration lies on general graphs and trees, with the exception of cycles in the time model \cite{Higashikawa2012} and $n\times n$-grid graphs \cite{Ortolf2012}.\\
The goal of this paper is to analyze exploration strategies for other restricted graph classes in terms of their competitive ratios, which describe the relationship between the exploration time
achieved by an online algorithm and the optimal result of an (offline) algorithm with full knowledge of a given graph. The analysis will be done for both, the time and the energy model. 
The particular graph classes considered are cycles (in the energy model) and tadpole graphs, which consist of a cycle and a path attached to one of its nodes.

\subsection{Contribution}

We show that the exploration strategy of avoiding the longest edge has a competitive ratio of $1.5$ on cycles in the energy model with two agents and a competitive ratio of at least $2$
on tadpole graphs in the time model, when using three agents. We then provide a modified exploration strategy, which does not decide based on the next edge weights, but on the length of the paths traversed
by the agents up to that point. On cycles this modified strategy still has a competitive ratio of $1.5$ in the time model, but also achieves a competitive ratio of $1$ in the energy model using two agents.
On tadpole graphs our exploration strategy achieves a competitive ratio of $1.5$ with a team of four agents in the time model, 
and a competitive ratio of $1$ with a team of three agents in the energy model. With a team of two agents, we achieve a competitive ratio
of $2.5$ on tadpole graphs in both models. An overview of the results is given in Table \ref{resultstable}.

\begin{table}[H]
    \small
    \begin{adjustbox}{center}
    \begin{tabular}{|l|ll|ll|}
    \hline
                                  & \multicolumn{2}{l|}{Time Model}                                     & \multicolumn{2}{l|}{Energy Model}              \\ \hline
    Graph Class (\# of Agents)                            & \multicolumn{1}{c|}{Lower Bound} & \multicolumn{1}{c|}{Upper Bound} & \multicolumn{1}{l|}{Lower Bound} & Upper Bound \\ \hline
    Cycles ($2$ Agents)             & \multicolumn{1}{l|}{$1.5$ \cite{Higashikawa2012}}         & $1.5$ \cite{Higashikawa2012}                              & \multicolumn{1}{l|}{$1$}           & $1$ (Thm. \ref{ampCycles})           \\ \hline
    Tadpole Graphs ($2$ Agents)     & \multicolumn{1}{l|}{$1.5$ (Thm. \ref{TadpoleTimeLower})}        & $2.5$ (Thm. \ref{TadpoleUpperTwo})                             & \multicolumn{1}{l|}{$1.5$ (Thm. \ref{TadpoleEnergyLower})}         & $2.5$ (Thm. \ref{TadpoleUpperTwo})      \\ \hline
    Tadpole Graphs ($3$ Agents)     & \multicolumn{1}{l|}{$1.5$ (Thm. \ref{TadpoleTimeLower})}          & $2$ (Thm. \ref{TadpoleUpperThree})                             & \multicolumn{1}{l|}{$1$}           & $1$ (Thm. \ref{TadpoleUpperThree})          \\ \hline
    Tadpole Graphs ($4+$ Agents)     & \multicolumn{1}{l|}{$1.5$ (Thm. \ref{TadpoleTimeLower})}       & $1.5$ (Thm. \ref{TadpoleUpperFour})                             & \multicolumn{1}{l|}{$1$}           & $1$ (Thm. \ref{TadpoleUpperThree})          \\ \hline
    \end{tabular}
    \end{adjustbox}
    \caption{The new lower and upper bounds for the exploration of cycles and tadpole graphs.}
    \label{resultstable}
\end{table}

\subsection{Related Work}

\paragraph*{Single-Agent Graph Exploration} Online Graph Exploration with a single agent has been studied extensively for many different graph classes: For planar
graphs Kalyanasundaram and Pruhs developed the $16$-competitive \textit{ShortCut} algorithm \cite{Kalyanasundaram1993},
which was generalized by Megow et al.\ to the \textit{Blocking} algorithm, which has a competitive ratio of $16(1 + 2g)$ for
any graph with genus $g$ \cite{Megow2012}.\\
Miyazaki et al.\ gave optimal algorithms for cycles and unit-weight graphs, with competitive ratios and lower bounds of
$1+\sqrt{3}$ and $2$ respectively. Megow et al.\ have extended the result on unit-weight graphs, showing that online exploration of graphs with $c$ 
distinct weights is $2c$ competitive \cite{Megow2012}. Brandt et al.\ have shown that greedy exploration is $2$-competitive
for the case of tadpole graphs \cite{Brandt2020}, Fritsch examined the \textit{Blocking} algorithm on unicyclic and cactus graphs
and proved that it is $3$-competitive for unicyclic and $\frac{5}{2}+\sqrt{2}$ competitive on cactus graphs \cite{Fritsch2021}. 
Foerster and Wattenhofer analyzed different variations of directed graphs and provided matching upper and lower bounds of $n-1$ for 
deterministic algorithms on general directed graphs, where the lower bound could be improved to $\frac{n}{4}$ when using randomized
algorithms \cite{Foerster2016}. Very recently Baligács et al.\ have shown that the \textit{Blocking} algorithm achieves constant competitive
ratio on minor-free graphs \cite{Baligacs2023}.\\
On general graphs the best known strategy is the \textit{Nearest Neighbor} algorithm, with a competitive ratio of $\Theta(\log(n))$ shown by Rosenkrantz et al.\ \cite{Rosenkrantz1977}.
The best known lower bound of $\frac{10}{3}$ was shown by Birx et al.\ and holds for general graphs and planar graphs as well. \cite{Birx2020}. An overview of the graph classes that have
already been investigated for single-agent graph exploration is provided in Table \ref{single-agent-table}.

\begin{table}[H]
    \small
    \centering
    \begin{tabular}{|l|l|l|}
    \hline
                       & \multicolumn{1}{c|}{Lower Bound}                                                  & \multicolumn{1}{c|}{Upper Bound} \\ \hline
    General Graphs     & $10/3$ \cite{Birx2020}                                                                               & $O(\log(n))$ \cite{Rosenkrantz1977}                           \\ \hline
    Planar Graphs     & $10/3$ \cite{Birx2020}                                                                              & $16$ \cite{Kalyanasundaram1993}                              \\ \hline
    Unit-weight Graphs  & $2$\cite{MIYAZAKI2009}                                                                                 & $2$\cite{MIYAZAKI2009}                                \\ \hline
    $c$ distinct weights  & $2$\cite{MIYAZAKI2009}                                                                                 & $2c$\cite{Megow2012}                                \\ \hline
    Cycles             &  $\frac{1 + \sqrt{3}}{2}$\cite{MIYAZAKI2009}                                                                           & $\frac{1 + \sqrt{3}}{2}$\cite{MIYAZAKI2009}                           \\ \hline
    Tadpole Graphs    & $2$ \cite{Brandt2020}                                                                                  & $2$ \cite{Brandt2020}                                 \\ \hline
    Unicyclic Graphs   & $2$ \cite{Brandt2020}                                                                                  & $3$ \cite{Fritsch2021}                                \\ \hline
    Cactus Graphs  & $2$ \cite{Brandt2020}                                                                                  & $2.5+\sqrt{2}$ \cite{Fritsch2021}                           \\ \hline
    Directed Graphs    & \begin{tabular}[c]{@{}l@{}}$n-1$ (deterministic)\\$n/4$ (randomized)\end{tabular} \cite{Foerster2016} & $n-1$\cite{Foerster2016}                              \\ \hline
    \end{tabular}
    \caption{Current results for single-agent graph exploration.}
    \label{single-agent-table}
    \end{table}

\paragraph*{Multi-Agent Graph Exploration on general graphs and trees} In the case of multi-agent graph exploration the main focus of the literature lies on trees. For $k$ agents, Fraigniaud et al.\ gave an $O(\frac{k}{\log k})$-competitive algorithm
for the exploration of trees \cite{Fraigniaud2006}. For the case of unrestricted communication between agents this was improved very recently by Cosson and Massouli{\'e} \cite{DBLP:conf/innovations/CossonM24}, who developed an $O(\sqrt{k})$ competitive tree exploration algorithm.
In the case when $k = 2^{\omega(\sqrt{\log D \log \log D})}$ and $n = 2^{O(2^{\sqrt{\log D}})}$ (where $D$ is the diameter of the tree and $n$ the number of nodes) Ortolf and Schindelhauer presented a subpolynomial $k^{o(1)}$-competitive
algorithm \cite{Ortolf2014}.
Fraigniaud et al.\ also provided a lower bound of $\Omega(2-\frac{1}{k})$ even for the case of global communication, which was later improved to $\Omega(\frac{\log k}{\log \log k})$ by Dynia et al.\ 
\cite{Dynia}. Dynia et al.\ \cite{Dynia2006} looked at a different cost model, in which not the general exploration time, but the distance traversed by any
individual agent has to be minimized. They developed a $4-\frac{2}{k}$-competitive algorithm for the exploration of trees 
 and gave a lower bound of $1.5$. Dynia et al.\ \cite{Dynia2006a} have analyzed the relationship between height and density of a tree and the competitive 
ratio of exploration strategies. They provided an algorithm with a competitive ratio of $O(D^{(1-\frac{1}{p})})$ where $p$ describes the density of a given tree. Dereniowski et al.\ gave an exploration strategy for general graphs using $Dn^{1+\Omega(1)}$ agents, that explores a graph in $O(D)$ time steps,
which implies a constant competitive ratio. \cite{Dereniowski2015}. This was extended on trees by Disser et al.\ to show a lower bound of $\omega(1)$ where $n \log^c n \leq k \leq Dn^{1+o(1)}$ \cite{Disser2020}.

\paragraph*{Multi-Agent Graph Exploration on restricted graph classes} Higashikawa et al.\ gave an optimal $1.5$-competitive algorithm for cycles using global communication and a lower bound of $\Omega(\frac{k}{\log k})$ for trees when using greedy algorithms, 
which implies that the $\frac{k}{\log k}$-competitive algorithm by Fraigniaud et al.\ cannot be improved by using greedy strategies \cite{Higashikawa2012}. 
Ortolf and Schindelhauer gave an upper bound of $O(\log^2 n)$ for $k$-agent exploration of $n \times n$ grid graphs with rectangular holes when using local communication
and a lower bound of $\Omega(\frac{\log k}{\log \log k})$ for deterministic as well as $\Omega(\sqrt{\frac{\log k}{\log \log k}})$ for randomized strategies \cite{Ortolf2012}.

\paragraph*{Other Graph Exploration variations} Aside from the previous examples, there exist different variations of multi-agent graph exploration, like exploration of unlabeled graphs using devices called pebbles \cite{Bender2002, Disser2018},
exploration using battery constrained agents \cite{Bampas2018}, exploration where edges are opaque \cite{Brass2011}, or exploration where nodes are space restricted and can only be occupied by one agent at any given time \cite{Czyzowicz2017}.
A $2$-competitive algorithm for exploration of cycles was developed by Osula for a model in which the exploration is started without any agents and spawning a new agent incurs an invocation cost $q$ \cite{Osula2017}.

\subsection{Outline}\label{outline}
In \S \ref{preliminaries} we define the exploration and cost models used. 
In addition, we describe the selected graph classes and some general concepts of online computation.\\
In \S \ref{ale_tad}, we investigate the competitiveness of the \textit{ALE algorithm} on tadpole graphs.
In \S \ref{cycles} to \ref{ntadpoles} we provide a modification of the \textit{ALE algorithm} and analyze its competitive ratio for the different graph classes, starting with cycles
in \S \ref{cycles}, followed by strategies for tadpole graphs using two to four agents in \S \ref{tadpoles}. Finally, a discussion of the results is given in \S \ref{discussion}.